\chapter{Introduction}
\label{chap:introduction}

This blueprint describes the mathematics formalized in the Lean library
\href{https://github.com/glaroton/lc-analysis}{lc-analysis}: the differential and integral
calculus of Michal and Bastiani on locally convex spaces. Every definition and theorem below is
linked to its Lean counterpart, and the dependency graph shows how the results build on each
other. All the statements in this document are proved in Lean, with no \texttt{sorry}, and the
main theorems depend only on the standard axioms of Lean.

\section*{Why Michal--Bastiani calculus}

Infinite-dimensional Lie groups and manifolds of mappings are modeled on spaces that are not
normable: Fréchet spaces of smooth functions such as $C^\infty(M)$, spaces of test functions,
direct limits. On such spaces the Fréchet derivative is not available, since there is no norm
with which to measure the remainder. One is faced with the problem of choosing a notion of
differentiability that does not depend on a norm, and that still has a chain rule, a mean value
theorem and a fundamental theorem of calculus.

The answer that has become standard in infinite-dimensional Lie theory goes back to Michal
(1938, \cite{Mic38}) and Bastiani (1964, \cite{Bas64}). It was developed further by Keller
\cite{Kel74}, Hamilton \cite{Ham82}, Milnor \cite{Mil84}, Glöckner \cite{Glo02} and Neeb
\cite{Nee06}, and its main reference today is the book of Glöckner and Neeb \cite{GN26}, which
we follow closely; \cite[1.2.13]{GN26} refers to its numbered result 1.2.13. We remark that
this is not the convenient setting of Kriegl and Michor \cite{KM97}: smoothness is defined here
through the topology of the model space, and not through smooth curves.

The idea is to ask for less than in the Banach case. A map is $\Ck{1}$ if all its directional
derivatives exist and depend continuously on the point \emph{and} on the direction, jointly. No
uniformity in the direction is required. In normed spaces this notion is weaker than Fréchet
differentiability, but the two give the same smooth maps (Chapter~\ref{chap:frechet}).

Mathlib \cite{mC20}, the mathematical library of the Lean proof assistant \cite{dMU21}, has a
complete Fréchet calculus on normed spaces, and a theory of topological vector spaces and
locally convex spaces. It does not have a differential calculus beyond the normed case. The
purpose of lc-analysis is to fill this gap, as a foundation for the formalization of
infinite-dimensional Lie theory.

\section*{Standing assumptions}

Unless otherwise stated, $E$, $F$ and $G$ denote real locally convex Hausdorff topological
vector spaces, $U \subseteq E$ is open, and $f \colon U \to F$. These spaces are not assumed to
be complete in any sense; whenever some form of completeness is needed (mainly for the
existence of integrals) it appears as an explicit hypothesis. The Lean statements sometimes
need less than the standing assumptions: for instance, the definition of weak integrals makes
sense in any topological vector space. The precise hypotheses can be read in the linked Lean
declarations.

Two tools replace norm estimates throughout:
\begin{itemize}
  \item \textbf{Hahn--Banach.} Continuous linear functionals separate points from closed convex
  sets. In particular, a vector is determined by its images under all continuous linear
  functionals.
  \item \textbf{Local convexity.} The closed convex neighborhoods of $0$ form a basis of the
  topology.
\end{itemize}
Most proofs reduce a statement in $F$ to statements about real functions, by composing with
continuous linear functionals, and then the one-variable calculus of Mathlib applies. Where in
the normed case a norm estimate would bound an increment, here we show instead that the
increment lies in a prescribed closed convex neighborhood.

\section*{Contents}

In Chapter~\ref{chap:differentiability} we define $\Ck{k}$ maps and prove the first-order
theory: linearity of the differential, the mean value theorem, the chain rule, partial
derivatives and the symmetry of the second differential. Chapter~\ref{chap:higher} compares
the recursive definition of $\Ck{k}$ with Bastiani's definition through iterated directional
derivatives, and Chapter~\ref{chap:rules} collects further rules of calculus (bilinear maps,
Taylor's formula, the second-order chain rule). Chapters~\ref{chap:integrals}
to~\ref{chap:mackey-lipschitz} are devoted to integration: weak integrals and their basic
properties, the completeness hypotheses under which they exist, and the equivalence of two
definitions of Mackey completeness. Finally, Chapter~\ref{chap:frechet} compares
Michal--Bastiani and Fréchet differentiability on normed spaces.

\section*{What is new and what is not}

The mathematics is classical: every result in this document appears in the cited literature,
and the proofs follow it. The few places where the argument differs from the sources are
pointed out in remarks. As far as we know, Michal--Bastiani calculus had not been formalized
before, neither in Lean nor in another proof assistant.

\section*{How the library was written}

@glaroton set the mathematical program, chose the framework and the references, fixed
the design constraints and reviewed the results. The Lean proofs were written by the AI
assistant Claude (Anthropic), working under his direction. Correctness does not rest on trusting
either of them: Lean's kernel checks every proof, and the statements are kept close to the
literature so that they can be compared with it directly.
